% ====================================================================
%  发明构建书（内部工作文档，非递交件）
%  由 notes.sty 排版。要改内容直接改本文件，这是唯一来源。
% ====================================================================
\documentclass[11pt,UTF8]{ctexart}
\usepackage{notes}
\notestitle{发明构建书}

\begin{document}

\nH{发明构建书（A1）}

{\ntbf 发明名称}：一种大语言模型智能体的工具调用控制方法及系统

{\ntbf 申请类型}：发明专利（不可走实用新型——实用新型仅保护产品的形状、构造或其结合，不接受方法权利要求）

{\ntbf 目标局}：CNIPA

{\ntbf IPC 建议}：G06F 9/48（进程调度/执行控制）、G06F 16/903（检索）、G06N 3/006 或 G06N 5/04（推理）

\par\vspace{2mm}\hrule\vspace{3mm}

\nHH{一、核心发明构思}

智能体运行时无法凭工具调用次数判断"是否在推进"。若干次措辞不同的调用可能反复取回同一个页面、同一条状态、同一个空结果或同一个失败。现有做法要么不判环，要么以"同一调用重复 k 次"这一单一条件判环——前者放任空转，后者在智能体场景里既误杀（同一入参在轮询中本就应重复）又漏判（措辞微变即绕过）。

本发明的核心构思是：{\ntbf 把工具观察的证据身份从具体的输出表示中解耦出来，使原始输出、外置引用、有限预览和压缩恢复表示在指向同一内容时保持同一证据身份；再把该身份、规范化规则版本和调用结果写入可重放的证据账本，以账本状态为主条件、以调用侧与结果侧的重复计数为合取副条件，在执行前阻断已被证明无效的重复调用。}

三个不可省略的落点：

\begin{enumerate}
  \item {\ntbf 内容摘要值是跨表示的证据身份}（桥接 A2）。大输出被外置为产物文件后，会话上下文里只剩引用式结果，其中的存储路径包含调用标识，因而同一份内容在不同调用下路径必不相同。若以路径、预览或调用标识作为证据身份，外置本身就会把"重复"伪装成"新增"，循环检测随外置一起失效。本发明以内容摘要值作为跨表示的证据身份依据，使外置、压缩恢复与循环检测共存。
  \item {\ntbf 三重合取阻断}。无新证据连续计数达阈值、同一调用指纹的调用计数达阈值、同一调用指纹的同结果计数达阈值，三者同时满足才拒绝。任一单独条件都会误杀合法行为。
  \item {\ntbf 账本可按版本确定性重放}。观察记录保存规范化规则版本、表示类型和完成顺序；会话上下文被压缩或恢复后，本发明从持久化观察记录或已存储的"工具调用—工具结果"配对按同一规则重放全部计数与集合，使压缩不导致"忘记自己在打转"，且重放不重新调用模型或外部工具。
\end{enumerate}

\nHH{二、要解决的技术问题}

\begin{enumerate}
  \item 智能体在一个回合内反复发起不产生新观察的工具调用，持续消耗算力、令牌与外部网络请求配额，且不产生任何可用信息增量；
  \item 以"调用是否重复"为唯一判据的抑制手段，对轮询/状态查询类工具误杀，对措辞微变的同义调用漏判；
  \item 工具输出中的时间戳、耗时、请求标识等易变字段使内容未变的结果每次都呈现为"新结果"，导致判重失效；
  \item 大体积工具输出被外置存储后，会话上下文中仅保留含调用标识的存储路径，使同一内容在判重时被反复认定为新证据；
  \item 上下文压缩后，前述判重所依赖的状态量随原始消息一并丢失，压缩后的会话无法察觉其在压缩前已陷入空转。
\end{enumerate}

\nHH{三、技术方案要素分解}

\begin{ntable}{Y{0.22}Y{0.22}Y{0.22}Y{0.22}}
{\ntbf 编号} & {\ntbf 技术特征} & {\ntbf 类别} & {\ntbf 落位} \\
\nthead
F1 & 对每次已完成调用分别生成调用指纹与证据指纹 & 核心 & 权1 \\
F2 & 调用指纹 = 工具名 + 规范化入参 & 核心 & 权1 \\
F3 & 证据指纹生成前做结构规范化（键名排序）并剔除易变字段 & 核心 & 权1 \\
F4 & 超阈值输出外置为产物文件 + 引用式结果替换 & 核心（A2 桥接） & 权1 \\
F5 & {\ntbf 以内容摘要值而非存储路径作为证据指纹依据} & {\ntbf 创造性落点①} & 权1 \\
F6 & 维护调用指纹计数、证据指纹集合、同结果计数、无新证据连续计数 & 核心 & 权1 \\
F7 & 新证据清零连续计数；重复或空则累加 & 核心 & 权1 \\
F8 & 达第一阈值 → 注入策略变更指示 & 核心 & 权1 \\
F9 & {\ntbf 达第二阈值 ∧ 调用计数≥第三阈值 ∧ 同结果计数≥第四阈值 → 执行前拒绝} & {\ntbf 创造性落点②} & 权1 \\
F10 & 轮询类工具第二阈值取整数倍 & 支撑 & 权1 \\
F11 & {\ntbf 压缩时从工具调用—工具结果配对重建全部状态量} & {\ntbf 创造性落点③} & 权1 \\
F12 & 错误输出用错误签名（剔除数字变量）作规范化结果 & 从属 & 权3 \\
F13 & 合成拒绝结果自身不计入账本、不清零连续计数 & 核心 & 权1 \\
F14 & 策略变更指示的去抖（同一连续状态内不重复注入） & 实施方式 & 说明书 \\
F15 & 关闭/仅提示/阻断三档运行模式 & 从属 & 权5 \\
F16 & 外置阈值由上下文预算推导并上下界钳制 & 从属（A2） & 权7 \\
F17 & 引用式结果携带首尾保留、中部省略的预览与检索提示 & 从属（A2） & 权8 \\
F18 & 产物文件以内容摘要值前缀命名，同内容幂等复用 & 从属（A2） & 权7 \\
F19 & 重建的回合边界取自最后一条非内部用户消息 & 从属 & 权9 \\
F20 & 账本聚合事实写入压缩摘要与固定上下文 & 从属 & 权10 \\
F21 & 并行调用组按已完成结果判定，不回溯撤销已派发调用 & 从属 & 权11 \\
F22 & 事件记录只输出元数据，不复制原始输出 & 从属 & 权12 \\
F23 & 系统/装置 & 独立 & 权13 \\
F24 & 计算机可读存储介质 & 独立 & 权14 \\
\end{ntable}

\nHH{四、有益效果（结构性推导，不含实验数据）}

\begin{itemize}
  \item 由于判据取自工具输出内容而非调用次数，措辞不同但取回同一观察的调用被识别为同一证据，空转得以在其发生时被察觉；
  \item 由于剔除易变字段后再生成证据指纹，内容未变的状态查询不会因时间推移而被误判为新证据；
  \item 由于以内容摘要值而非含调用标识的存储路径作为证据身份，大输出外置不再破坏循环检测，二者可同时启用；
  \item 由于阻断需三重条件合取且轮询类工具阈值放大，正当的重试与轮询不被误杀，而已被证明无效的重复调用在执行前即被拒绝，从而不再产生对应的算力、令牌与外部网络请求消耗；
  \item 由于状态量可从消息历史重建，上下文压缩不导致判重能力丢失，长回合中的空转同样可被察觉。
\end{itemize}

\nHH{五、不写入申请文件的内容（纪律）}

\begin{itemize}
  \item 不写入任何实验数据、命中率、节省比例、对比性能；
  \item 不写入具体的商业产品名、仓库名、环境变量名以外的实现标识；
  \item 不在背景技术中承认现有技术的优越性，不主张未经代理人核实的文献号；
  \item 术语一次锁定，全文不换词（见下）。
\end{itemize}

\nHH{六、术语表（全文锁定）}

\begin{ntable}{Y{0.3}Y{0.3}Y{0.3}}
{\ntbf 术语} & {\ntbf 含义} & {\ntbf 禁止替换为} \\
\nthead
智能体运行时 & 承载模型—工具迭代循环的执行主体 & 框架、引擎、agent \\
用户回合 & 从一条用户消息开始到本轮作答结束的过程 & 会话、轮次、turn \\
工具调用 / 工具结果 & 模型发起的调用 / 执行后返回的输出 & 函数调用、action \\
调用指纹 & 由工具名与规范化入参生成的摘要值 & 输入哈希 \\
证据指纹 & 由规范化后的工具输出生成的摘要值 & 输出哈希 \\
证据指纹集合 & 本回合已出现过的证据指纹的集合 & 去重表 \\
调用指纹计数 & 同一调用指纹出现的次数 & 重复次数 \\
同结果计数 & 同一调用指纹连续产生同一结果的次数 & —— \\
无新证据连续计数 & 连续多少次已完成调用未产生新证据 & streak \\
易变字段 & 时间戳、耗时、请求标识、链路追踪标识、心跳等字段 & 噪声字段 \\
内容摘要值 & 对工具输出原文计算的摘要值 & 哈希、sha256 \\
产物文件 & 存放超阈值工具输出原文的本地文件 & 缓存文件 \\
引用式结果 & 替换进会话上下文的、含路径与摘要值的结构化结果 & 占位符 \\
策略变更指示 & 注入工具结果中的、要求改变假设/来源/目标/机制的内部提示 & 提醒、nudge \\
轮询类工具 & 用于等待或查询状态、其重复调用属正当行为的工具 & 状态工具 \\
上下文压缩 & 为释放上下文预算而对历史消息做的摘要化处理 & 总结、compaction \\
\end{ntable}

\end{document}
